\documentclass[UTF8]{ctexart}

%\RequirePackage{amsmath} \RequirePackage{amssymb}
\usepackage{amscd,amsthm,amsfonts,amssymb,amsmath}
\usepackage[colorlinks=true]{hyperref}
%\usepackage{fullpage}
\usepackage[all]{xy}
\usepackage{mathrsfs}
\usepackage{tikz,mathpazo}
%\usetikzlibrary{quotes,angles}
%\usetikzlibrary{calc}
%\usetikzlibrary{decorations.pathreplacing}
\usepackage{bm}
\usepackage{geometry}
\usepackage{xcolor}
\usepackage{eso-pic}
\usepackage{CJK}
\usepackage{arydshln}
\usepackage{mathptmx}
\usepackage[11pt]{moresize}



\newcommand{\BC}{{\mathbb {C}}}

\newcommand{\BN}{{\mathbb {N}}}
\newcommand{\BQ}{{\mathbb {Q}}}
\newcommand{\BR}{{\mathbb {R}}}
\newcommand{\BZ}{{\mathbb {Z}}}
\newcommand{\bv}{{\bm v}}
\newcommand{\bw}{{\bm w}}
\newcommand{\bu}{{\bm u}}
\newcommand{\bx}{{\bm x}}
\newcommand{\by}{{\bm y}}

\newcommand{\bb}{{\bm b}}
\newcommand{\Coker}{{\mathrm{Coker}}}



\newcommand{\End}{{\mathrm{End}}}
\newcommand{\GL}{{\mathrm{GL}}}

\newcommand{\SO}{{\mathrm{GSO}}}
\newcommand{\Hom}{{\mathrm{Hom}}}

\renewcommand{\Im}{{\mathrm{Im}}}

\newcommand{\Isom}{{\mathrm{Isom}}}
\newcommand{\Id}{{\mathrm{Id}}}



\newcommand{\Ker}{{\mathrm{Ker}}}



\newcommand{\rank}{{\mathrm{rank}}}


\newcommand{\SU}{{\mathrm{SU}}}
\newcommand{\Sym}{{\mathrm{Sym}}}

\newcommand{\sub}{{\mathrm{sub}}}

\newcommand{\tr}{{\mathrm{tr}}}

\newcommand{\matrixx}[4]{\begin{pmatrix}
		#1 & #2 \\ #3 & #4
	\end{pmatrix} }
	
	
	
	\newcommand{\wt}{\widetilde}
	\newcommand{\wh}{\widehat}
	\newcommand{\pp}{\frac{\partial\bar\partial}{\pi i}}
	\newcommand{\pair}[1]{\langle {#1} \rangle}
	\newcommand{\wpair}[1]{\left\{{#1}\right\}}
	\newcommand{\intn}[1]{\left( {#1} \right)}
	\newcommand{\norm}[1]{\|{#1}\|}
	\newcommand{\sfrac}[2]{\left( \frac {#1}{#2}\right)}
	\newcommand{\ds}{\displaystyle}
	\newcommand{\ov}{\overline}
	\newcommand{\incl}{\hookrightarrow}
	\newcommand{\lra}{\longrightarrow}
	\newcommand{\lla}{\longleftarrow}
	\newcommand{\ra}{\rightarrow}
	\newcommand{\imp}{\Longrightarrow}
	\newcommand{\lto}{\longmapsto}
	\newcommand{\bs}{\backslash}
	
	
	
	
	\newtheorem{thm}{定理}
	\newtheorem{cor}[thm]{推论}
	\newtheorem{lem}[thm]{引理}
	\newtheorem{prop}[thm]{命题}
	\newtheorem{defn}[thm]{定义}
	\newtheorem{ques}[thm]{问题}
	\newtheorem{eg}[thm]{例题}
	\newtheorem{ex}[thm]{习题}
	\newtheorem{rmk}[thm]{注释}
		\newtheorem{ntn}[thm]{记号}

	\newcommand{\sk}{\medskip}\newcommand{\bsk}{\bigskip}
	\newcommand{\s}{\sk\noindent}\newcommand{\bigs}{\bsk\noindent}
	
	
	%accented words

	\def\mat(#1,#2,#3,#4){
		\left[\begin{array}{cc}
			#1 & #2 \\ #3 & #4\end{array}
		\right]
	}
		\def\clm(#1,#2){
		\left[\begin{array}{c}
			#1 \\ #2 \end{array}
\right]
	}
		\def\clmn(#1,#2){
		\left[\begin{array}{c}
			#1 \\ \vdots\\ #2 \end{array}
\right]
	}
		\def\clmt(#1,#2,#3){
		\left[\begin{array}{c}
			#1 \\ #2  \\ #3\end{array}
\right]
	}
		\def\clmf(#1,#2,#3){
		\left[\begin{array}{c}
			#1 \\ #2  \\ \vdots \\ #3\end{array}
\right]
	}


		
\begin{document}

\title{习题课材料（二）简要解答}

\date{}

\maketitle

\vspace{-1cm}




\noindent{\Large\textbf{注: 带 $\heartsuit  $号的习题有一定的难度、比较耗时, 请量力为之.}}





\begin{ex}[矩阵的对角化]
	计算：
	\begin{enumerate}
		\item$
		\begin{bmatrix}
		2&3\\5&7\\
		\end{bmatrix}
		\begin{bmatrix}
		2&0\\0&3\\
		\end{bmatrix}
		\begin{bmatrix}
		-7 & 3\\5&-2
		\end{bmatrix}$;
		\item$
		\begin{bmatrix}
		-7 & 3\\5&-2
		\end{bmatrix}
		\begin{bmatrix}
		2&3\\5&7\\
		\end{bmatrix}$;
		\item$
		\begin{bmatrix}
		17&-6\\35&-12\\
		\end{bmatrix}^5$.
		\item$
		\begin{bmatrix}
		1&3&1\\2&2&1\\3&4&2\\
		\end{bmatrix}
		\begin{bmatrix}
		1&&\\&2&\\&&1\\
		\end{bmatrix}
		\begin{bmatrix}
		0&2&-1\\1&1&-1\\-2&-5&4\\
		\end{bmatrix}$;
		\item$
		\begin{bmatrix}
		8&6&-6\\4&6&-4\\8&8&-6\\
		\end{bmatrix}^6$.
	\end{enumerate}
\end{ex}

\begin{ex}
	对 $n$ 阶方阵 $A$, 设 ${\rm Com}(A)=\{ n\text{ 阶方阵 }B: AB=BA\}$. ({\rm Com}表示Commutator.)
	\begin{enumerate}	
		\item 证明: 任取 $B, C \in {\rm Com}(A)$, 都有 $I_n, kB+\ell C, BC \in {\rm Com}(A)$, 其中$k,\ell \in \mathbb{R}$.
		\item 设 $A=\begin{bmatrix}
		1& 1& 0\\
		0& 1& 1\\
		0& 0& 1\\
		\end{bmatrix}$. 证明：$J_3=\begin{bmatrix}
		0& 1& 0\\
		0& 0& 1\\
		0& 0& 0\\
		\end{bmatrix} \in {\rm Com}(A)$, 而且 ${\rm Com}(A)=\{k_1I_3+k_2J_3+k_3J_3^2:k_1, k_2, k_3 \in \mathbb{R}\}$.
	\end{enumerate}		
\end{ex}



\begin{ex}
	是否存在矩阵 $A$ 满足：存在矩阵 $X$ 使得 $XA=I$, 但不存在 $Y$ 使得 $AY=I$? 有没有方阵满足上述条件？
\end{ex}




\begin{ex}
	定义函数 $\tr: M_n(\mathbb{R}) \to \mathbb{R}, A \mapsto \tr(A):=a_{11}+ \cdots + a_{nn}$, 为取方阵的对角线元素之和. $\tr(A)$ 称为方阵 $A$ 的迹.
	\begin{enumerate}
		\item 证明 $\tr$ 满足如下三个条件：
		\begin{itemize}
			\item $\tr(kA+\ell B)=k\tr(A)+\ell\tr(B)$, $k,\ell \in \mathbb{R}.$
			\item $\tr(AB)=\tr(BA)$；
			\item $\tr(I_n)=n$.
		\end{itemize}
		\item 说明是否存在 $A, B$, 使得 $AB-BA=I_n$.
		 \item ($\heartsuit$)  在$n=2$时证明满足上述三个条件的函数一定是 $\tr$. 		\item ($\heartsuit$) 对一般的$n$证明满足上述三个条件的函数一定是 $\tr$.
	\end{enumerate}	
\end{ex}


\begin{proof}[提示]
我们只给出4的提示：这里我们用点与之前不同的记号: $M_{ij}$ 指第 $(i,j)$ 项等于 $1$, 其它所有项等于 $0$ 的矩阵; $P_{ij}$ 是把单位矩阵 $I$ 的 $(i,i)$ 项和 $(j,j)$ 项改成零, 再把 $(i,j)$ 项和 $(j,i)$ 项改成 $1$, 这样得到的矩阵. 直接计算得 $$P_{ij}M_{ii}P_{ij}^{-1}=M_{jj},\quad M_{ij}M_{k\ell}=\left\{\begin{array}{ll} 0  & j\neq k \\ M_{i\ell} & j=k\end{array}\right.$$
	
	首先由第一个条件, 令 $A$ 为零矩阵 ${\bm 0}_n$ 知 $\tr(\bm{0}_n)=0$. 由第二个条件知
	$$\tr(M_{jj})=\tr(P_{ij}M_{ii}P_{ij}^{-1})=\tr(P_{ij}M_{ii}P_{ij}^{-1}P_{ij})=\tr(M_{ii}),$$
	而由第一条
	$$\tr(I_n)=\tr(M_{11}+\cdots+M_{nn})=\tr(M_{11})+\cdots+\tr(M_{nn});$$
	由第三条可知 $\tr(M_{ii})=1$. 另一方面, 当 $i\neq j$ 时
	$$\tr(M_{ij})=\tr(M_{i1}M_{1j})=\tr(M_{1j}M_{i1})=\tr({\bm 0}_n)=0.$$
	于是
	$$	\tr(A)=\sum_{i,j}a_{ij}\tr(M_{ij})=a_{11}+\cdots+a_{nn}.$$
\end{proof}

\begin{ex}\label{ex-14}
	\begin{enumerate}
		\item 任取 $m \times n$ 矩阵 $X$; 证明: 分块矩阵
		$\left( \begin{array}{cc}
		I_m & X \\
		0 & I_n
		\end{array} \right)$
		可逆, 并求其逆矩阵.
		\item 对分块矩阵 $\left( \begin{array}{cc}
		A & B \\
		C & D
		\end{array} \right)$, 计算 $\left( \begin{array}{cc}
		I_m & X \\
		0 & I_n
		\end{array} \right)\left( \begin{array}{cc}
		A & B \\
		C & D
		\end{array} \right)$. 由此判断矩阵
		$$	\begin{bmatrix}
		1& 0& 0&\cdots&0&c_1\\
		0& 1& 0&\cdots&0&c_2\\
		0& 0& 1&\cdots&0&c_3\\
		\vdots& \vdots& \vdots&\ddots&\vdots&\vdots\\
		0 & 0 & 0 & \cdots & 1 & c_{n-1}\\
		b_1& b_2& b_3&\cdots& b_{n-1}&a\\
		\end{bmatrix}$$ 何时可逆, 并在它可逆时计算它的逆.
		\item 设 $A \in M_m(\mathbb R)$, $B \in M_{m \times n}(\mathbb R)$,
		$C \in M_n(\mathbb R)$,	证明：若 $A, C$ 可逆,	则分块矩阵
		$\left( \begin{array}{cc}
		A & B \\
		0 & C
		\end{array} \right)$
		也可逆, 并求其逆矩阵.
	\end{enumerate}
\end{ex}
\begin{proof}[提示]
	2. 把给定的矩阵分块: $\mat(I_{n-1}, \alpha,\beta^T,a)$, 作分块初等行列变换得到
	$\mat(I_{n-1},0,0,a-\beta^T\alpha)$.
\end{proof}

\begin{ex}
	\begin{enumerate}
		\item 对 $n$ 阶可逆矩阵 $A$ 和 $n$ 维列向量 $\bu, \bv$, 设 $1+\bv^TA^{-1}\bu \ne 0$, 证明：$A+\bu\bv^T$ 可逆, 且
		$$(A+\bu\bv^T)^{-1}=A^{-1}-\dfrac{1}{1+\bv^TA^{-1}\bu}A^{-1}\bu\bv^TA^{-1}.$$ (这称为Sherman-Morrison-Woodbury公式.)
		
		\item 设 $a_i > 0(1 \le i\le n),$ 求矩阵：
		$$\begin{bmatrix}
		a_1+1& 1& 1&\cdots&1\\
		1& a_2+1& 1&\cdots&1\\
		1& 1& a_3+1&\cdots&1\\
		\vdots& \vdots& \vdots&\cdots&\vdots\\
		1& 1& 1&\cdots& a_n+1\\
		\end{bmatrix}$$
		的逆矩阵.
	\end{enumerate}	
\end{ex}
\begin{proof}[提示]
	当然, 我们可以做形式计算: 将等式右边右乘上 $A+\bu\bv^T$. 我们也可以借助习题\ref{ex-14}, 令 $\alpha=A^{-1}\bu,\beta=-\bv$ 可知矩阵 $\mat(I_n, A^{-1}\bu,-\bv^T, 1)$ 可逆. 作 (分块) 初等行列变换:
	$$\mat(I_n, A^{-1}\bu,-\bv^T, 1)\to \mat(A,\bu,-\bv^T,1)\to\mat(A+\bu\bv^T,\bu,0,1).$$
	考虑这个变换的过程, 以及矩阵 $\mat(I_n, A^{-1}\bu,-\bv^T, 1)$ 的逆矩阵, 即\ref{ex-14}题中的公式.
	
	再令 $A={\rm diag}(a_1,\cdots,a_n)$, $\bu=\bv=\clmn(1,1)$ 可得 2.
\end{proof}




\begin{ex}
	如果矩阵 $A=\begin{bmatrix}
	a_{ij}
	\end{bmatrix}_{n\times n}$ 满足: 对于任意 $i=1,\dots,n$, 都有 $|a_{ii}|>\sum_{j\ne i}|a_{ij}|$, 则称其为(行)对角占优的. 证明：对角占优矩阵一定可逆.
\end{ex}
\begin{proof}[提示]
	我们证明 $A\bx=0$ 只有零解. 反证: 设 $\bx\neq0$ 是它的解, 取 $\ell$ 使得 $|x_\ell|=\max_{1\leq i\leq n}|x_i|>0$. 由 $\sum_{j=1}^na_{\ell j}x_j=0$ 知
	$$|a_{\ell\ell}x_\ell|=\left|\sum_{j\neq \ell}a_{\ell j}x_j\right|\leq\sum_{j\neq\ell}|a_{\ell j}||x_j|\leq \sum_{j\neq\ell}|a_{\ell j}||x_\ell|<|a_{\ell\ell}||x_\ell|,$$
	矛盾.
\end{proof}
\begin{ex}
	证明：
	\begin{enumerate}
		\item 任意方阵 $A$ 都可以唯一地表为 $A=B+C$, 其中 $B$ 是对称矩阵, $C$ 是反对称矩阵．
		\item ($\heartsuit$) $n$ 阶方阵 $A$ 是反对称矩阵当且仅当任取 $n$ 维列向量 $\bx$, 都有 $\bx^{T}A\bx=0$.
		\item 设 $A, B$ 是对称矩阵, 则$A=B$ 当且仅当任取 $n$ 维列向量 $\bx$, 都有 $\bx^{T}A\bx=\bx^{T}B\bx$.
	\end{enumerate}
\end{ex}
\begin{proof}[提示]
	1. $A=\frac{A+A^T}{2}+\frac{A-A^T}{2}$.
	
	2. 必要性: 若 $A$ 反对称, 则 $(\bx^TA\bx)^T=\bx^TA^T\bx=-\bx^TA\bx$. 充分性: 若对任意 $\bx$, $\bx^TA\bx=0$, 令 $\bx=e_i+e_j$ 知 $a_{ij}=-a_{ji}$.
	
	3. 证明充分性: 由于对任意 $\bx$ 有 $\bx^T(A-B)\bx=0$, 由 2 知 $A-B$ 是反对称矩阵. 已知 $A-B$ 是对称矩阵, 所以 $A-B=0$.
\end{proof}


\begin{ex}[$\heartsuit$]
	$A$ 为 $n$ 阶实方阵. 证明以下结论：
	\begin{enumerate}
		\item 若对于任意的 $n$ 维实列向量 $\bx$, 都有 $(A\bx)\cdot (A\bx)=\bx\cdot\bx$, 则 $A$ 必须是正交矩阵．
		\item 若对于任意两个 $n$ 维实列向量 $\bx,\by$, 都有 $(A\bx)\cdot \by=\bx\cdot(A\by)$, 则 $A$ 必须是对称矩阵．
	\end{enumerate}
\end{ex}
\begin{proof}[提示]
	1. 令 $\bx=\bu+\bv$ 有 $(A\bu)\cdot(A\bv)=\bu\cdot\bv$. 取 $\bu=e_i,\bv=e_j$, 因为 $Ae_i$ 是 $A$ 的第 $i$ 个列向量 $\bx_i$, 我们有 $\bx_i\cdot\bx_j=\delta_{ij}$, 故 $A^TA=I_n$.
	
	2. 取 $\bx=e_i, \by=e_j$, 有 $a_{ji}=(Ae_i)\cdot e_j=e_i\cdot(Ae_j)=a_{ij}$.
\end{proof}



\noindent{\itshape 思考题}:
对于$n$阶方阵$A$试说明下列几条为何等价：
\begin{enumerate}
		\item $A$可逆．
		\item 任取$n$维向量$\bb$，方程组$A \bx = \bb$ 的解唯一，且解为$\bx = A^{-1}\bb$.
\item 齐次方程组$A\bx = {\bf 0}$只有零解．
\item $A$对应的阶梯型矩阵有$n$个主元.
\item $A$对应的行简化阶梯形矩阵一定是$I_n$.
\item $A$是有限个初等矩阵的乘积．
	\end{enumerate}



\end{document}
